% Online Supplement S.C — full grids behind the validation audits of
% submitted Appendix C (sec_app03). Permutation draws, rolling-origin
% tables, environment-dominance grids, granularity sweeps, timing-set
% register, leakage cross-tabulation, and the direct-defendant timing grid.
\section{Validation Audits: Full Grids}
\label{supp:validation}

This section holds the complete grids behind the validation audits of
the submitted Appendix~C. Definitions, headline numbers, the
over-crediting proposition, and the interpretation each audit licenses
are in that appendix; here we report the underlying draws and
specifications. ``CADE-defendant contact'' is an adjudication-anchored
exposure label throughout: a firm shares a tender-item with a directly
adjudicated CADE defendant, which records proximity, not liability.

\subsection{Opportunity-Cell Support and Expected Contact}
\label{supp:opp_cells}

The three nested opportunity-cell definitions (COARSE $=$ item group
$\times$ year; MEDIUM $=$ item group $\times$ year $\times$ buyer; STRICT
$=$ item code $\times$ year $\times$ buyer) trade resolution against
support. The mean cell defendant-contact rate $\bar p_g$ is
$0.017$--$0.019$ across definitions. Under STRICT, $99.5\%$ of cells hold
fewer than five tender-items and $75.7\%$ are singletons; the strict
common-support rule retains about $46\%$ of firms under MEDIUM cells and
about $9\%$ under STRICT. Mean excess $X_i = O_i - E_i$ concentrates at
high participation tiers ($+0.068$ at $T\ge 50$ versus $-0.006$ to
$-0.000$ at lower tiers), but the level of $O_i$ is dominated by $E_i$:
exposure does the heavy lifting and excess is the residual the score
must rank. The full cell-construction, expected-contact, and
common-support retention grids are produced by the opportunity-cell
build referenced in the submitted appendix.

\subsection{Companion Audits to the Opportunity-Adjusted Battery}
\label{supp:perm_matched}

Three companion exercises accompany the control-function ladder of
Table~\ref{tab:control_function_submission} (submitted Appendix~C) and
do not corroborate a robust residual signal net of opportunity.
\emph{(i)} A matched-stratum label permutation places the observed
PR-AUC at $p=0.127$ (not significant). \emph{(ii)} A firm-exposure
simulation yields a higher level PR-AUC than the raw score ($p=1.000$).
\emph{(iii)} A CEM opportunity-matched comparison gives a within-stratum
AUC of $0.626$ ($n=6{,}039$, full common support), and the within-stratum
change in $\Pr(\text{cobidder})$ for FL14 versus non-FL is in fact
\emph{negative} ($-0.017$). Under a Benjamini--Hochberg correction over
the full correlated family (nested increment, within-stratum AUC,
CEM-matched AUC, two exposure-preserving permutations, and the
intensity-restricted re-run), the only test that survives at the $5\%$
level is the $+0.010$ nested DeLong increment ($p=0.013$); the headline
within-stratum and permutation nulls are non-rejections to begin with.

\paragraph{Intensity-restricted label sensitivity.} Restricting
positives to cobidders sharing at least two distinct defendant
tender-items ($368$ of $651$ positives) sharpens the negative
conclusion. Raw discrimination rises mechanically (ROC-AUC $0.854$), but
the exposure-only model rises faster ($0.873$ on the exposed subsample;
$0.956$ unconditionally); the within-opportunity-stratum AUC remains at
chance ($0.506$); the nested increment falls to $+0.003$ ($p=0.47$); and
the matched-stratum permutation again does not reject ($p=0.39$).
% src: outputs/sensitivity_contact2/ (02b_opportunity_sensitivity_contact2.R, 2026-06-04)

\subsection{Audit of the Audit: Leakage Tiers, Falsifiability, Power}
\label{supp:audit_armor}

% src: outputs/diagnostics/audit_armor/ (scripts 12/12b, 2026-06-04)
\emph{Exposure-benchmark leakage tiers.} Contact-derived benchmarks
partially encode a contact-defined label. Ranking by observed contact
$O_i$ gives AUC $0.905$ and a plug-in $E_i$ $0.985$ (pure label
encoding); the firm-leave-one-out $E_i$ used in the main text gives
$\valArmorExpLOO$; recomputing every cell's defendant density from
non-cobidder rows only---so no eventual positive contributes to any
firm's expected contact---leaves a label-blind opportunity ranking of
$\valArmorExpLB$. The exposure benchmark is therefore an upper bound on
opportunity's role.

\emph{Falsifiability of the within-stratum null.} Within the same
opportunity-decile strata that put the score at chance, a
positive-control score ($O_i$) ranks the label at
$\valArmorOiControlMedium$: the design detects a residual when one
exists. The headline within-stratum AUC ($\valExpWithinAUC$) stratifies
on opportunity deciles; rebuilding the strata from the discrete
expected-contact cells leaves the within-stratum score AUC in the range
$\valArmorWithinMedium$--$\valArmorWithinStrict$ across granularities
(coarse $\valArmorWithinCoarse$, medium $\valArmorWithinMedium$, strict
$\valArmorWithinStrict$). The sequence is non-monotone---it dips at the
medium definition and rises in the strictest cells---and retained
support shrinks from $15{,}246$ to $615$ firms across that range. The
strict-cell $\valArmorWithinStrict$ is not evidence of a residual: it
rests on only $615$ retained firms ($16{,}295$ within-cell pairs), where
the $O_i$ positive control itself falls to $0.740$, signalling that the
strictest cells are too thin to discriminate reliably. Stratifying
instead on the \emph{label-blind} expected contact leaves the score at
$\valArmorWithinLBall$; the negative controls in the submitted
Appendix~D identify that ordering as generic co-participation geometry,
not cartel-specific signal.

\emph{Power of the permutation design.} Injecting synthetic
within-stratum signal into the labels and re-running the matched-strata
permutation ($B=200$, $60$ simulations per effect size) shows the test
is correctly sized under the null (rejection rate $0.05$ at zero
injected signal) and rejects with probability $\valArmorPowerTwo$ at a
true within-stratum AUC of $0.52$, $\valArmorPowerFive$ at $0.55$, and
$\valArmorPowerTen$ at $0.60$. The observed non-rejection therefore
bounds any surviving residual below a within-stratum AUC of roughly
$0.55$.

\emph{Label-frozen timing.} Freezing the pool, the score, and the label
ingredients to $2009$--$2016$ (always-loser status from training-window
wins only; defendant contact from training-window tender-items only)
leaves $\valArmorFrozenPool$ rankable incumbents. The frozen score ranks
new $2017$--$2019$ defendant contact at AUC $\valArmorFrozenProspAUC$
($\valArmorFrozenProspN$ positives) and fully frozen retrospective
contact at $\valArmorFrozenRetroAUC$ ($\valArmorFrozenRetroN$
positives)---participation volume forecasting generic future contact
with active firms, carrying no cartel-specific content. These
label-frozen counts are distinct objects from the strict-holdout rows in
Table~\ref{tab:strict_holdout_submission}: the $\valArmorFrozenPool$
pool comprises all frozen-always-loser incumbents with any pre-2017
history, whereas the ``training-AL \& test-active'' row intersects that
pool with test-window activity ($11{,}351$ firms, $312$ positives under
the full-window label). $\valArmorFrozenProspN$ counts firms whose
\emph{first} defendant contact arises in $2017$--$2019$, while
$\valArmorFrozenRetroN$ counts contact within $2009$--$2016$ only, and
the two need not nest the full-window label counts.

\subsection{Timing Information Sets and Rolling-Origin Validation}
\label{supp:timing_sets}

The CADE files carry only judgment dates (\texttt{data\_julgamento},
mostly 2015--2025) and \emph{no conduct dates}, so all timing is
\emph{year-level} (BEC participation year), never day-level
(\texttt{DAY\_LEVEL\_TIMING\_UNAVAILABLE}). We define seven candidate
information sets, A--G, distinguished by which inputs (score, zero-win
status, threshold, label) are drawn pre- versus post-window:

\begin{itemize}
\item[\textbf{A}] Labeled retrospective upper bound (all inputs
full-window).
\item[\textbf{B}] Leak-flagged decomposition (zero-win status from full
period; flagged as leaking future wins).
\item[\textbf{C/D/G}] Honest incumbent-triage designs: strict
prospective score, rankable incumbent pool, retrospective label, read as
triage rather than prospective deployment. We lead on these.
\item[\textbf{E}] Discloses the entrant coverage hole (entrants carry
score zero with no training history).
\item[\textbf{F}] Adjudication-observable-at-time: infeasible because all
judgment dates post-date a 2009--2016 freeze.
\end{itemize}

The strict holdout (submitted Appendix~C,
Table~\ref{tab:strict_holdout_submission}) freezes the score,
always-loser status, and median$+1.5\times$IQR threshold on 2009--2016
(train-window threshold $\valStrictThreshTrain$, vs.\ full-window
$13.5$) and evaluates against 2017--2019 cobidder labels.

\paragraph{Rolling-origin validation.} Expanding the training window for
$t=2014,\dots,2019$ posts PR-AUC $\approx 0.012$--$0.014$ and
precision@$500=$recall@$500=0$ on the full universe in every origin
year, with the worst ROC below chance ($\valRollWorstAUC$, 2015). The
within-pool (training-always-loser) ROC stays in the $0.66$--$0.70$
range across origins, again an artifact of a thicker retrospective
training history rather than real-time learning.

\paragraph{Leakage decomposition.} The in-sample item-level reference
AUC ($\valLeakRefAUC$) decomposes so that the surviving structural
component stays in the $\valLeakStructLow$--$\valLeakStructHigh$ range
under cobidder-firm holdout (out-of-fold AUC $\valLeakCVAUC$) and
temporal holdout (AUC $\valAUCFLfirmTemp$), with the
$\valLeakTautLow$--$\valLeakTautHigh$ AUC lost to mechanical reuse
disclosed as such.

\subsection{Case Composition: Environment-Dominance Grids}
\label{supp:environment_dominance}

The operational signal concentrates by case, item group, and year.
Dropping the single largest case (\texttt{trens\_metros}) lowers PR-AUC
from $\valLOCOfullPRAUC$ to $\valLOCOdropLargestPRAUC$ ($-37\%$);
dropping the top two cases takes PR-AUC to $0.058$; case-balanced
precision@$500$ (weighting each case equally) is $\valCaseBalancedPrec$
against the firm-weighted pooled $\valFullPrec$. The largest case holds
$\valTopCasePosShare$ of positives and $\valTopCaseTP$ of the true
positives in the top $500$. The largest item group holds
$\valTopItemGroupShare$ of positives (item-group HHI $0.100$) while the
buyer distribution is diffuse (buyer HHI $0.009$); dropping the largest
item group leaves PR-AUC essentially unchanged ($0.143 \to 0.134$).

A clustered randomization-inference test permuting cobidder labels
within item-group $\times$ year clusters ($B=1{,}000$) clears the null
on the ranking metrics (ROC, PR-AUC, precision@$500$ all
$p=\valClusterRIp$), but top-$500$ case coverage is \emph{not}
significant ($p=\valClusterRIcovP$), consistent with one case anchoring
the operational signal. The case-composition inference rests on a small
number of adjudicated cartels---twelve matched CADE cases, of which six
are informative for the case-grouped exercises---so the effective number
of clusters at the natural unit of treatment (the case) is small. The
randomization inference is clustered at the finer item-group $\times$
year level; under the honest few-effective-clusters reading,
$p=\valClusterRIp$ and the coverage $p=\valClusterRIcovP$ are
directionally informative rather than exactly size-controlled.

\subsection{Direct-Defendant Timing Grid}
\label{supp:direct_defendant_timing}

The FL-binary score (retrospective incumbent triage) ranks direct CADE
defendants at ROC $\approx \valDirectStrictAUC$ under both full and
strict timing ($0.491$ and $0.492$), i.e.\ at chance, with
precision@$500=0$. Continuous participation volume carries some
information about defendants---ROC $0.66$--$0.70$, modestly above
chance---but with precision@$500\le0.004$, so it does not operationally
retrieve them either. The asymmetry is the design: the loser-side
\emph{binary} flag is uninformative about defendants, who typically sit
on the winner side of the arrangement, while participation volume is not
entirely silent.
